%==============================================================================
% BODY — ported from papers/study_05/MANUSCRIPT_DRAFT_v2.md (Codex-gate-passed,
% CODEX-S5-MANUSCRIPT-GATE-20260910-001, verdict PASS_TO_TYPESETTING_WITH_LIMITATIONS).
% No class-specific commands here: paste into the UEH template unchanged.
% Headline numbers come from % GENERATED by code/study_05_crosscountry/tables/build_tables.py -- DO NOT EDIT BY HAND
% In-text statistics. Regenerate via build_tables.py; never edit by hand.
\newcommand{\NAdults}{100,560}
\newcommand{\NEconomies}{97}
\newcommand{\NMechAdults}{90,118}
\newcommand{\NMechEconomies}{89}
\newcommand{\FormalBorrowMean}{9.2}
\newcommand{\MOneBeta}{3.617}
\newcommand{\MOneSE}{0.514}
\newcommand{\MOneCIlo}{2.597}
\newcommand{\MOneCIhi}{4.638}
\newcommand{\MTwoBetaOne}{3.957}
\newcommand{\MTwoBetaThree}{-1.975}
\newcommand{\MTwoBetaThreeSE}{0.474}
\newcommand{\MTwoBetaThreeCIlo}{-2.916}
\newcommand{\MTwoBetaThreeCIhi}{-1.033}
\newcommand{\MTwoBetaThreeP}{$<$0.0001}
\newcommand{\MTwoBmedianBeta}{-3.958}
\newcommand{\MTwoBmedianSE}{0.970}
\newcommand{\MTwoBmedianP}{$<$0.0001}
\newcommand{\MtwocF}{6.77}
\newcommand{\MtwocP}{0.002}
\newcommand{\MThreeAReceive}{4.080}
\newcommand{\MThreeAMake}{5.172}
\newcommand{\MThreeBReceive}{-0.269}
\newcommand{\MThreeBReceiveP}{0.587}
\newcommand{\MThreeBMake}{-2.021}
\newcommand{\MThreeBMakeP}{$<$0.0001}
\newcommand{\InformalBetaThree}{1.736}
\newcommand{\InformalBetaThreeP}{0.005}
\newcommand{\InformalBetaThreeQ}{0.010}
\newcommand{\InformalN}{100,428}
\newcommand{\CtrlxModBeta}{-1.545}
\newcommand{\CtrlxModSE}{0.451}
\newcommand{\CtrlxModP}{0.00091}
\newcommand{\CtrlxModCIlo}{-2.44}
\newcommand{\CtrlxModCIhi}{-0.65}
\newcommand{\WithinRtwoBase}{0.0481}
\newcommand{\WithinRtwoFull}{0.0491}
\newcommand{\DBIntBeta}{-1.955}
\newcommand{\DBIntP}{0.00010}
\newcommand{\DBIntN}{95,909}
\newcommand{\DBIntG}{95}
\newcommand{\MDEpp}{1.33}
\newcommand{\VnmSlope}{5.07}
\newcommand{\VnmSlopeCIlo}{3.87}
\newcommand{\VnmSlopeCIhi}{6.27}
\newcommand{\VnmCoverage}{59.4}
\newcommand{\SampleMedianCoverage}{40.4}
\newcommand{\ThinCoverageEconomies}{24}
\newcommand{\ThinShareN}{23,458}
\newcommand{\ThinSharePct}{23.3}
\newcommand{\RawRateNoDig}{3.78}
\newcommand{\RawRateDig}{14.13}
\newcommand{\RawGap}{10.3}
\newcommand{\RegionSSA}{34}
\newcommand{\RegionECA}{17}
\newcommand{\RegionLAC}{16}
\newcommand{\RegionMENA}{10}
\newcommand{\RegionEAP}{9}
\newcommand{\RegionSA}{5}
\newcommand{\RegionHIC}{6}
; tables from \input{tables/*}.
% The eight residual limitations from the Codex review are preserved verbatim in
% Section 7 ("Residual limitations that must survive typesetting").
%==============================================================================

\begin{abstract}
\noindent We ask whether digital-payment activity is associated with formal borrowing
among adults, and whether that association is larger where credit-information coverage is
thinner --- the individual-level analogue of the firm-level pattern documented by
\citet{GalileaFaraziMare2026}. Using the 2024 Global Findex microdata for \NAdults{} adults
in \NEconomies{} economies, we estimate survey-weighted linear probability models with
economy fixed effects and economy-clustered standard errors, following a pre-registered
plan. Making or receiving a digital payment is associated with a \MOneBeta{}
percentage-point-higher probability of having borrowed from a formal financial institution
in the past 12 months (95\% CI \MOneCIlo{}--\MOneCIhi{} pp), conditional on account
ownership, demographics and economy fixed effects. The pre-registered moderation prediction
is \textbf{not supported}: the interaction of digital-payment activity with \emph{lower}
2019 credit-information coverage is negative (\MTwoBetaThree{} pp per economy-level standard
deviation; 95\% CI \MTwoBetaThreeCIlo{} to \MTwoBetaThreeCIhi{} pp), the opposite of the
hypothesised sign. For the registered coverage moderator, the negative interaction persists
across weighting, sample and control-set changes and all \NEconomies{} leave-one-economy-out
fits. The observed clustered precision gives an approximate 80\%-power minimum detectable
interaction of \MDEpp{} pp per SD. Two alternative institutional proxies are imprecise and
do not reproduce a significant reversal. The registered receive-versus-make mechanism and
the informal-borrowing placebo also fail. We interpret the results as a cautious, non-causal
external-population check that does not reproduce the credit-information-substitution reading
at the individual level, and we discuss why a coverage floor and the account-access channel
are plausible explanations. We report Vietnam's position in the cross-country distribution
for local context.

\medskip
\noindent\textbf{Keywords:} digital payments, financial inclusion, formal borrowing, credit
information, Global Findex, pre-registration. \quad\textbf{JEL:} G21, G51, O16, D14.
\end{abstract}

\section{Introduction}

Digital payments have become widely used in low- and middle-income economies: the 2024
Global Findex shows that a majority of adults make or receive at least one digital payment
\citep{KlapperSingerStaritaNorris2025}. A recurring policy claim is that the transaction
record such payments generate can serve as an alternative credit signal, extending formal
credit to borrowers whom traditional credit-information systems do not cover.
\citet{GalileaFaraziMare2026} provide the most direct evidence for this idea: across 101
economies, firms that use electronic payments report weaker credit constraints, and the
association is stronger where credit-information depth is lower --- consistent with
electronic-payment records substituting for a missing credit file.

Whether the same pattern generalises to \emph{individuals} remains an open empirical
question. Our targeted, non-systematic search found no prior multi-country individual-level
test of the exact Findex interaction, but cannot establish priority. Individuals face
different lenders, different products and different information technologies than firms, and
the Global Findex --- the natural data source at the individual level --- has been used
mainly to describe digital-payment adoption and broad financial-inclusion outcomes rather
than to test this specific moderation. \citet{BergBurgGombovicPuri2020} show that
digital-footprint variables can predict individual default and add information to a
credit-bureau score. Their complementarity result is a competing explanation, not evidence
for our pre-registered substitution prediction, and their setting is a single e-commerce
lender in one country. \citet{BazarbashBeaton2020} relate cross-country fintech credit to
financial development and information infrastructure using provider-level data, again not at
the individual level.

This paper runs a pre-registered individual-level test. We do \textbf{not} claim novelty of
the question or of the method; we position the exercise as an individual-level descriptive
analogue and an external-population check of the established firm-level pattern. Three
features discipline the exercise. First, the full analysis plan --- sample, variables,
specifications, robustness checks, multiple-testing correction and decision rules --- was
frozen and hash-committed before any outcome model was estimated. Second, the primary
results were produced by one implementation and then independently reproduced by a separate
cold implementation that did not see the first; 41 of 44 compared quantities match to better
than $5\times10^{-13}$. Third, the decision rules were written to make a null or contrary
interaction a valid, reportable outcome rather than a failure.

The pre-registration matters because the headline result is a contrary one. We find that
digital-payment activity is positively associated with formal borrowing within economies
(\MOneBeta{} pp, 95\% CI \MOneCIlo{}--\MOneCIhi{}), but that this association is
\textbf{smaller}, not larger, in economies where 2019 credit-bureau and credit-registry
coverage was lower (interaction \MTwoBetaThree{} pp per standard deviation, 95\% CI
\MTwoBetaThreeCIlo{} to \MTwoBetaThreeCIhi{}). The sign reversal is robust to unweighted and
population weights, to excluding high-income economies, to adding country-level macro
controls and their interactions, to dropping zero-coverage economies, to leaving out each of
seven world regions, and to leaving out each of the \NEconomies{} economies one at a time.
The depth and legal-rights proxies are statistically imprecise and do not establish a
general institutional gradient. The registered mechanism split (receiving versus making
payments) and the informal-borrowing placebo also do not behave as the substitution reading
predicts.

Our contribution is therefore an honest, reproducible negative result. At the individual
level, in this cross-section, the cross-country pattern that motivates the
digital-footprint-substitution policy narrative does not appear; if anything the data show
the opposite gradient. We evaluate plausible explanations --- a coverage floor in the
thinnest-coverage economies, and the possibility that in those economies digital-payment
activity mainly proxies for basic account access rather than for an incremental credit
signal --- and we are explicit about what the design cannot rule out.

\section{Related literature and positioning}

\paragraph{Firm-level evidence.} \citet{GalileaFaraziMare2026} is the closest paper. Using
World Bank Enterprise Survey data for 101 economies, they document that electronic-payment
use is associated with weaker firm credit constraints and that the association is larger
where the depth-of-credit-information index is lower. Our design mirrors theirs at the
individual level: same broad hypothesis, same moderation logic, different unit and different
data. Their robustness section changes firm-level outcomes, payment instruments, fixed
effects and controls within the Enterprise Surveys; neither the paper nor its robustness
appendix uses Global Findex or estimates an individual-level borrowing specification. Thus
it sharpens, rather than duplicates, the external-population test conducted here.

\paragraph{Digital footprints and credit.} \citet{BergBurgGombovicPuri2020} show that easily
accessible digital-footprint variables match or exceed a credit bureau score in predicting
consumer default at a German e-commerce lender and, crucially, complement rather than
replace bureau information. This establishes that the \emph{information content} of a
digital trace can be real; it does not establish that the trace expands formal borrowing at
the population level, still less that the effect concentrates where bureaus are thin.

\paragraph{Cross-country digital credit.} \citet{BazarbashBeaton2020} find a market-level
gap-filling pattern, while \citet{CornelliEtAl2023} conclude that fintech and big-tech
credit complement rather than substitute for bank credit. Both study aggregate provider
volumes rather than an individual borrower's use of formal credit. The disagreement suggests
that substitution and complementarity can depend on the level of analysis and the lending
channel.

\paragraph{Closest individual-level evidence.} \citet{AlokGhoshKulkarniPuri2024} study
India's UPI rollout with a causal design and find credit growth concentrated in places where
account access was more recent; only fintech lenders expanded lending to new-to-credit
borrowers. Their within-country setting, treatment and outcomes differ sharply from ours,
but the result makes basic formal-system access a serious competing explanation for the
cross-country gradient.

\paragraph{Global Findex.} \citet{KlapperSingerStaritaNorris2025} and the associated
database document the rapid adoption of digital payments and its correlation with account
ownership and with saving and borrowing behaviour. The database is descriptive by
construction. We use it for a specific pre-registered moderation test rather than for
adoption description.

\paragraph{Positioning.} We claim only an individual-level descriptive analogue and
external-population check of the \citeauthor{GalileaFaraziMare2026} pattern. A targeted
search found no prior individual-level, multi-country Findex study of this exact
interaction, but search coverage is incomplete and does not justify a categorical priority
claim. We make no causal or mechanism-identification claim; novelty remains a candidate
distinction rather than an established fact.

\section{Data}

\subsection{Individual data}

The individual data are the 2024 wave of the World Bank Global Findex microdata
(\url{https://microdata.worldbank.org/catalog/7860}; DOI 10.48529/bk9n-8r43; SHA-256
\texttt{ca307a0c\ldots f90fce}). The borrowing and digital-payment module was administered
in 98 of the 140 released economies; the reduced questionnaire elsewhere omits the outcome.
Our unit of observation is an adult aged 15 or older within an economy.

\begin{itemize}
\item \textbf{Outcome} \texttt{formal\_borrow}: one if the respondent reports having
borrowed from a bank or another type of formal financial institution in the past 12 months
($\mathit{fin22a}=1$), zero if not ($\mathit{fin22a}=2$); ``don't know'' and ``refused'' are
dropped.
\item \textbf{Exposure} \texttt{anydigpayment}: one if the respondent made or received at
least one digital payment in the past 12 months, zero otherwise.
\item \textbf{Mechanism split} (secondary): \texttt{receive\_dig}, one if the respondent
received any of wages, government transfers, a pension or an agricultural payment into an
account; \texttt{make\_dig}, one if the respondent made a digital merchant payment or paid a
utility bill from an account. Unknown responses are treated as missing, never as zero, and
are dropped only for the mechanism regressions.
\item \textbf{Controls} (frozen primary set): account ownership at a financial institution,
an indicator for female, age centred and its square, three education categories, five
within-economy income quintiles, and an urban indicator.
\end{itemize}

\subsection{Country moderator}

The moderator is the country's \textbf{2019 credit-information coverage}: the maximum of the
2019 private credit-bureau coverage and the 2019 public credit-registry coverage, in percent
of adults, from the archived World Bank Doing Business series
(\texttt{IC.CRED.ACC.PRVT.CRD.ZS}, \texttt{IC.CRED.ACC.PUBL.CRD.REG.COVR.ZS}; snapshot
SHA-256 \texttt{ef0c67f9\ldots bd06c}). We denote its negative $z$-score, computed across
the estimation economies, \lowcovz, so that higher values mean thinner 2019 coverage.
Across the \NEconomies{} estimation economies the raw coverage measure has a cross-economy
(unweighted) mean of 41.5\%, a median of \SampleMedianCoverage\%, and a wide spread
(\ThinCoverageEconomies{} economies below 10\%). The depth-of-credit-information index and
the strength-of-legal-rights index, both 2019 vintage, are used as distinct robustness
proxies, not as the same construct.

The moderator is measured five years before the outcome. Credit infrastructure is
slow-moving, so this is a defensible historical measure, but we treat the vintage mismatch
as a limitation and never describe the moderator as the contemporaneous 2024 system.

The source also requires an integrity qualification. The World Bank paused and later
discontinued \emph{Doing Business} after reviews of irregular data changes in the 2018 and
2020 report cycles. The Bank's December 2020 review specifically identified \emph{Getting
Credit} irregularities for China (DB2018) and Saudi Arabia (DB2020), both present in our
estimation sample. Our moderator uses archived 2019-vintage coverage fields rather than an
overall rank, so the review does not by itself establish that every value used here is
erroneous. We nevertheless treat the history as a measurement-risk warning and report a
sensitivity that excludes both economies.

As a partial external check, the 2019 bureau- and registry-coverage, depth and legal-rights
values were cross-validated against the published \emph{Doing Business 2020} economy
profiles for eight contrasting economies (Vietnam, Afghanistan, Kenya, Georgia, China, South
Africa, Saudi Arabia and Panama), and Vietnam's 2019 bureau coverage was additionally
confirmed against a live World Bank API query; all matched the archived snapshot, including
the extreme values (zero-depth Afghanistan, full-coverage Georgia and China). This does not
rule out error in the economies not checked, and the broad measurement-risk limitation
stands.

\subsection{Country controls (robustness only)}

For one robustness specification we add, at the economy level, log GDP per capita and
domestic credit to the private sector (\% of GDP), each the latest World Bank observation no
later than 2023, and the in-sample economy account-ownership rate, together with their
interactions with digital-payment activity. GDP per capita is 2023 for every economy; the
private-credit series has heterogeneous vintages (2008--2023), which we disclose.

\subsection{Sample construction}

Table~\ref{tab:sample} records the deterministic sample. Requiring a valid formal-borrowing
response leaves 102{,}644 adults in 98 economies. Requiring complete primary covariates
removes a further 2{,}084 adults and one economy (Lesotho, whose income-quintile variable is
entirely missing). The moderator merges to every remaining economy by ISO3 code. The
estimation sample is \textbf{\NAdults{} adults in \NEconomies{} economies}; the mechanism
sub-sample, which additionally requires non-missing receive and make indicators, is
\NMechAdults{} adults in \NMechEconomies{} economies. Vietnam contributes 998 respondents.

% GENERATED by code/study_05_crosscountry/tables/build_tables.py -- DO NOT EDIT BY HAND
\begin{table}[htbp]
\centering
\caption{Sample construction.}
\label{tab:sample}

\begin{adjustbox}{max width=\textwidth}
\begin{tabular}{l r r l}
\toprule
Step & {Adults} & {Economies} & \\
\midrule
Valid formal-borrowing response ($\mathit{fin22a}\in\{1,2\}$) & 102,644 & 98 & (42 economies exit) \\
\quad + non-missing digital-payment indicator & 102,644 & 98 &  \\
\quad + non-missing account ownership & 102,644 & 98 &  \\
\quad + non-missing sex & 102,644 & 98 &  \\
\quad + non-missing age & 102,644 & 98 &  \\
\quad + non-missing education & 102,245 & 98 &  \\
\quad + non-missing income quintile & 101,247 & 97 & (LSO exits) \\
\quad + non-missing urbanicity & 100,560 & 97 &  \\
\quad + 2019 coverage moderator merged by ISO3 & 100,560 & 97 &  \\
\bottomrule
\end{tabular}
\end{adjustbox}
\smallskip
\begin{minipage}{\linewidth}\footnotesize \emph{Notes.}
Deterministic listwise construction from the 2024 Global Findex microdata
(catalogue 7860). ``Don't know'' and ``refused'' on the outcome are dropped before
step 1. The moderator merges to every remaining economy.
\end{minipage}

\end{table}


Raw rates (survey-weighted): \RawRateNoDig\% of adults with no digital payment have borrowed
formally, against \RawRateDig\% of adults with a digital payment --- an unadjusted gap of
about \RawGap{} percentage points, which the regressions below reduce to \MOneBeta{} pp once
account ownership, demographics and economy fixed effects are held fixed.

The \NEconomies{}-economy frame spans all seven World Bank regions represented by the
module: \RegionSSA{} economies are in Sub-Saharan Africa, \RegionECA{} in Europe and Central
Asia, \RegionLAC{} in Latin America and the Caribbean, \RegionMENA{} in the Middle East and
North Africa, \RegionEAP{} in East Asia and the Pacific, \RegionSA{} in South Asia, and
\RegionHIC{} are classified as high income. The \ThinCoverageEconomies{} economies with
coverage below 10\% contribute \ThinShareN{} respondents, or \ThinSharePct\% of the
estimation rows. These are unweighted sample counts; the primary estimator gives every
economy equal total weight.

% GENERATED by code/study_05_crosscountry/tables/build_tables.py -- DO NOT EDIT BY HAND
\begin{table}[htbp]
\centering
\caption{Descriptive statistics.}
\label{tab:descriptives}

\begin{adjustbox}{max width=\textwidth}
\begin{tabular}{l r r r l}
\toprule
\multicolumn{5}{l}{\emph{Panel A. Individual characteristics (survey-weighted means)}} \\
 & {All} & {No digital} & {Digital} & \\
 & {($N=100,560$)} & {payment ($N=44,368$)} & {payment ($N=56,192$)} & \\
\midrule
Formal borrowing & 0.092 & 0.038 & 0.141 & \\
Account at a financial institution & 0.505 & 0.182 & 0.805 & \\
Female & 0.512 & 0.558 & 0.469 & \\
Age (years) & 38.6 & 38.2 & 39.1 & \\
Urban & 0.390 & 0.339 & 0.438 & \\
\midrule
\multicolumn{5}{l}{\emph{Panel B. Economy-level 2019 moderator ($G=97$ economies, unweighted)}} \\
 & {Mean} & {Median} & {[p25, p75]} & {Vietnam} \\
\cmidrule(lr){1-5}
2019 credit-information coverage (\%) & 41.5 & 40.4 & {[10.4, 59.9]} & 59.4 \\
\quad 2019 private-bureau coverage (\%) & 32.2 & 20.6 & {[0.0, 56.5]} & 20.6 \\
\quad 2019 public-registry coverage (\%) & 16.0 & 2.1 & {[0.0, 26.1]} & 59.4 \\
Depth-of-credit-information index (0--8) & 5.9 & 7.0 & {[6.0, 8.0]} & 8.0 \\
Strength-of-legal-rights index (0--12) & 6.0 & 6.0 & {[3.0, 8.0]} & 8.0 \\
\bottomrule
\end{tabular}
\end{adjustbox}
\smallskip
\begin{minipage}{\linewidth}\footnotesize \emph{Notes.}
Panel A: means weighted by \texttt{w\_equal} (each economy sums to one).
Panel B: one row per economy, unweighted; the moderator used in estimation is the
negative $z$-score of the coverage measure (\texttt{lowcov2019\_z}).
\end{minipage}

\end{table}


\begin{figure}[htbp]
\centering
\includegraphics[width=0.8\linewidth]{F1_country_gap}
\caption{Economy-level raw (survey-weighted) gap in formal borrowing between adults with and
without a digital payment, against 2019 credit-information coverage. One point per economy;
the fitted line is a simple OLS through the economy points; Vietnam is marked. Source:
2024 Global Findex; \NEconomies{} economies.}
\label{fig:country_gap}
\end{figure}

\section{Empirical strategy}

All specifications were frozen in the pre-registration and hash-committed before estimation.

\subsection{Primary specification}

For adult $i$ in economy $c$, let $y_{ic}$ be formal borrowing, $D_{ic}$ digital-payment
activity, $A_{ic}$ account ownership, and $L_c$ the lower-2019-coverage moderator
(\texttt{lowcov2019\_z}). We estimate, by weighted least squares,
\begin{align}
y_{ic} ={}& \beta_1 D_{ic} + \beta_2 A_{ic} + X_{ic}'\gamma + \delta_c + \varepsilon_{ic}
  \tag{M1}
\end{align}
and, for the moderation test,
\begin{align}
y_{ic} ={}& \beta_1 D_{ic} + \beta_3\,(D_{ic}\times L_c)
  + \beta_2 A_{ic} + \beta_4\,(A_{ic}\times L_c)
  + X_{ic}'\gamma + \delta_c + \varepsilon_{ic}.
  \tag{M2}
\end{align}
Here $X$ is the frozen control set, $\delta_c$ are economy fixed effects (which absorb the
moderator's main effect), and weights \texttt{w\_equal} make each economy sum to one, so the
estimand is the average within-economy association with economies weighted equally. Standard
errors are clustered by economy (\NEconomies{} clusters); we also report Rademacher
wild-cluster-bootstrap $p$-values (999 replications). \textbf{$\beta_3$ is the pre-registered
test.} $\beta_3>0$ would mean the digital-payment/formal-borrowing association is stronger
where 2019 coverage was thinner.

\subsection{Pre-registered decision rule}

The frozen rule: H-000500 is \emph{supported} only if $\beta_3>0$, significant at 5\%, and
robust to the below-median moderator contrast, the bureau/registry joint test, alternative
weighting, dropping high-income economies, the depth and legal-rights proxies, and
country-level controls, and if the mechanism split shows the receive interaction positive.
$\beta_3\approx0$ or $\beta_3<0$ means \emph{not supported}, to be reported plainly. A
non-null placebo interaction for informal borrowing also undermines the substitution
reading.

\subsection{Secondary and robustness specifications}

The plan specifies: a below-median moderator contrast (M2b); a bureau/registry joint test
(M2c); the receive-versus-make mechanism split (M3); three secondary outcomes with
Benjamini--Hochberg control over a frozen nine-hypothesis family (M4); four heterogeneity
interactions (M5); a ten-branch robustness grid (M6, including unweighted and population
weights, dummy-variable logit/probit, dropping high-income economies, the depth and
legal-rights proxies, country controls, a minimum economy-size filter, an internet-use
control, an economy-block bootstrap, a leave-one-economy-out sweep and an employment
control); and a descriptive Vietnam-positioning exercise (M7). Adversarial checks (drop each
control, restrict to account holders, drop zero-coverage economies, drop the lowest coverage
quartile, leave out each region, and a moderator sign-reversal algebra check) were run after
estimation.

\subsection{Interpretation limits}

The design is a single cross-section. Economy fixed effects remove all fixed cross-country
confounding but do not identify the effect of digital-payment use or of the moderator, and
do not make the cross-level interaction causal. Digital-payment activity and formal
borrowing are mechanically linked through account ownership, which we condition on but which
is itself a possible common consequence. We therefore use associational language throughout
and prohibit causal, policy-effect and mechanism-identification claims.

\section{Results}

\subsection{The association}

Digital-payment activity is associated with a \textbf{\MOneBeta{} percentage-point} higher
probability of formal borrowing (M1; economy-clustered SE \MOneSE{} pp;
$p\approx 3\times10^{-10}$; 95\% CI \MOneCIlo{}--\MOneCIhi{} pp; wild-cluster-bootstrap
$p=0.001$; \NAdults{} adults, \NEconomies{} economies). Against a sample formal-borrowing
rate of about \FormalBorrowMean\%, this is an association of roughly 40\% of the mean. It is
a within-economy, covariate-adjusted association, not a causal effect.
Table~\ref{tab:primary} reports M1 and M2 in full.

% GENERATED by code/study_05_crosscountry/tables/build_tables.py -- DO NOT EDIT BY HAND
\begin{table}[htbp]
\centering
\caption{Digital payments, formal borrowing and the 2019 coverage moderator.}
\label{tab:primary}

\begin{adjustbox}{max width=\textwidth}
\begin{tabular}{l r r}
\toprule
 & (1) M1 & (2) M2 \\
 & {Association} & {Moderation} \\
\midrule
Digital payment ($\beta_1$) & 3.617*** & 3.957*** \\
 & (0.514) & (0.482) \\[2pt]
Digital payment $\times$ lower 2019 coverage ($\beta_3$) & {--} & -1.975*** \\
 & {--} & (0.474) \\[2pt]
Account at FI ($\beta_2$) & 8.182*** & 7.800*** \\
 & (0.723) & (0.659) \\[2pt]
Account at FI $\times$ lower 2019 coverage ($\beta_4$) & {--} & 1.828*** \\
 & {--} & (0.616) \\[2pt]
\midrule
Demographic controls & {Yes} & {Yes} \\
Economy fixed effects & {Yes} & {Yes} \\
Clusters (economies) & {97} & {97} \\
Observations & {100,560} & {100,560} \\
Within-$R^2$ & {0.0475} & {0.0481} \\
\bottomrule
\end{tabular}
\end{adjustbox}
\smallskip
\begin{minipage}{\linewidth}\footnotesize \emph{Notes.}
Survey-weighted (\texttt{w\_equal}) linear probability models; coefficients and
standard errors in percentage points. Economy-clustered (CR1) standard errors in
parentheses. Demographic controls: female, age (centred) and its square, education
(3 categories), within-economy income quintile (5), urban. $\beta_3$ is the
pre-registered test. Rademacher wild-cluster-bootstrap $p$ (999 replications):
$\beta_1$ in M1, $p=0.001$; $\beta_3$ in M2, $p=0.001$.
Stars: *** $p<0.01$, ** $p<0.05$, * $p<0.10$ (economy-clustered).
\end{minipage}

\end{table}


\subsection{The pre-registered moderation test (H-000500)}

The moderation prediction is \textbf{not supported, with the interaction carrying the
opposite sign}. In M2 the interaction of digital-payment activity with lower 2019
credit-information coverage is \textbf{\MTwoBetaThree{} percentage points} per one
economy-level standard deviation (economy-clustered SE \MTwoBetaThreeSE{} pp;
$p\approx 6.8\times10^{-5}$; 95\% CI \MTwoBetaThreeCIlo{} to \MTwoBetaThreeCIhi{} pp;
wild-cluster-bootstrap $p=0.001$). The digital-payment/formal-borrowing association is
therefore \emph{weaker}, not stronger, in economies with thinner 2019 coverage --- the
reverse of the credit-information-substitution prediction. The below-median contrast (M2b)
agrees: the association is \MTwoBmedianBeta{} pp smaller in economies below the coverage
median (SE \MTwoBmedianSE{} pp; $p\approx 9.3\times10^{-5}$).

The observed clustered precision implies an approximate conventional 80\%-power minimum
detectable interaction of \MDEpp{} pp per SD at a two-sided 5\% level, and the 95\% interval
is entirely negative, so within this specification the data reject not only the registered
positive prediction but any non-negative interaction. This is a model-based precision
statement, not evidence that the true structural effect is negative.

Under the frozen decision rule this places H-000500 in the \emph{not supported} branch. We
do not relabel the reverse pattern as a confirmation of a different hypothesis; it was not
pre-registered.

\subsection{Robustness of the registered coverage result}

For the registered coverage moderator, the negative interaction persists across the
pre-registered weighting, sample and control-set branches and the adversarial
specifications (Table~\ref{tab:robustness}). The alternative depth and legal-rights proxies
do not yield statistically significant interactions and therefore do not generalise the
reversal to other institutional constructs.

% GENERATED by code/study_05_crosscountry/tables/build_tables.py -- DO NOT EDIT BY HAND
\begin{table}[htbp]
\centering
\caption{The M2 interaction $\beta_3$ across specifications (percentage points per SD).}
\label{tab:robustness}

\begin{adjustbox}{max width=\textwidth}
\begin{tabular}{l r l r r l}
\toprule
Specification & $\beta_3$ & $p$ & $N$ & $G$ & Note \\
\midrule
M2 baseline & -1.975 & $<$0.0001 & 100,560 & 97 & \\
\addlinespace
Unweighted & -2.129 & $<$0.0001 & 100,560 & 97 &  \\
Population weights & -2.783 & 0.00046 & 100,560 & 97 &  \\
Drop 6 high-income economies & -2.050 & $<$0.0001 & 94,682 & 91 & $z$ re-computed \\
+ country controls $\times$ digital payment & -2.354 & 0.003 & 100,560 & 97 &  \\
+ internet-use control & -2.048 & $<$0.0001 & 100,560 & 97 &  \\
+ employment control & -1.896 & 0.00017 & 96,925 & 96 &  \\
Economies with $N\geq 500$ & -1.975 & $<$0.0001 & 100,560 & 97 &  \\
\addlinespace
Moderator = depth-of-information proxy & -0.525 & 0.182 & 100,560 & 97 & same sign, imprecise \\
Moderator = legal-rights proxy & 0.705 & 0.236 & 100,560 & 97 & wrong sign, imprecise \\
\addlinespace
Restrict to account holders & -0.803 & 0.218 & 55,429 & 97 & imprecise \\
Drop zero-coverage economies & -2.121 & $<$0.0001 & 99,059 & 95 &  \\
Drop lowest coverage quartile & -2.880 & $<$0.0001 & 76,108 & 72 &  \\
Drop China and Saudi Arabia; re-standardize & -1.955 & 0.00010 & 95,909 & 95 & indep.\ reproduced \\
+ control $\times$ moderator (post-registration) & -1.545 & 0.00091 & 100,560 & 97 & post-reg., dual impl. \\
\addlinespace
Leave out each of 7 world regions & {-1.412 \ldots\ -2.187} & \multicolumn{3}{l}{all negative} & \\
Leave out each of 97 economies & {-1.737 \ldots\ -2.164} & \multicolumn{3}{l}{all negative; 0 sign reversals} & \\
Economy-block bootstrap (999) & -1.975 & \multicolumn{3}{l}{95\% percentile CI [-2.91, -1.05]} & \\
Moderator sign reversed (algebra check) & {1.975} & \multicolumn{3}{l}{exact negation of the baseline} & \\
\bottomrule
\end{tabular}
\end{adjustbox}
\smallskip
\begin{minipage}{\linewidth}\footnotesize \emph{Notes.}
All specifications are M2-form (percentage points per one economy-level SD of the
negative coverage $z$-score), economy-clustered (CR1). All rows were reproduced by the
independent cold implementation \texttt{CLAUDE-S5R-20260910-001}; the China/Saudi-Arabia
row was additionally reproduced by \texttt{CLAUDE-S5R-DB-INTEGRITY-20260910-001}
($|\Delta|=5\times10^{-14}$) and the control-by-moderator row by
\texttt{CODEX-S5-CTRLXMOD-REPRO-20260910-001} ($|\Delta|=3\times10^{-13}$). The
control-by-moderator row is a post-registration check outside the frozen quantity set.
Dummy-variable logit and probit interaction average marginal effects are
$-1.348$ and $-1.358$ percentage points (same sign).
\end{minipage}

\end{table}


\begin{figure}[htbp]
\centering
\includegraphics[width=0.8\linewidth]{F2_leave_one_out}
\caption{Leave-one-economy-out estimates of the M2 interaction $\beta_3$. Each point drops a
single economy and re-estimates; the dashed line is the full-sample estimate and the grey
line marks zero. All \NEconomies{} fits are negative.}
\label{fig:loo}
\end{figure}

Dummy-variable logit and probit give negative interaction average marginal effects
(approximately $-1.35$ and $-1.36$ pp from the primary run; the independent reproduction
obtains $-1.50$ and $-1.47$ pp, the same sign with a larger magnitude); the sign, not the
precise magnitude, is the informative quantity from these non-linear models. The
specifications that do not reach significance are those that also lose precision by
construction: the account-holders-only restriction and the two alternative moderator
constructs. Of the latter, the depth-of-credit-information proxy keeps the negative sign
($-0.525$ pp); only the legal-rights proxy turns positive ($+0.705$ pp), and it is far from
significant ($p\approx 0.24$) and measures strength of legal rights rather than coverage. No
coverage-based specification reverses toward the hypothesised sign.

One threat the frozen plan does not address is composition: M2 interacts only
digital-payment activity and account ownership with the moderator, so if the covariate mix
of digital-payment users varies with country coverage the interaction could be an artifact.
As a post-registration check, outside the frozen quantity set, we add each control
interacted with \lowcovz{} to M2. The interaction attenuates to \CtrlxModBeta{} pp per SD
(economy-clustered SE \CtrlxModSE{} pp; $p\approx\CtrlxModP$; 95\% CI \CtrlxModCIlo{} to
\CtrlxModCIhi{}) --- about a fifth of the baseline magnitude is absorbed once the control
slopes are allowed to vary with coverage, consistent with a compositional contribution ---
but it keeps its sign and significance. The reversal is therefore not a compositional
artifact of the control set, though composition plausibly accounts for part of the
gradient. A separate implementation reproduces this result within numerical tolerance. The
within-$R^2$ rises only from \WithinRtwoBase{} to \WithinRtwoFull{} (a single-implementation
diagnostic whose baseline value matches the frozen Stage-2 run exactly); we report this
transparent coefficient-stability comparison instead of treating a scalar selection bound as
identification for a cross-level interaction in a binary-outcome fixed-effects model.

\subsection{Mechanism and placebo}

The registered receive-versus-make mechanism (M3) does not behave as the substitution
reading predicts (Table~\ref{tab:mechanism}). In M3a the making-payment coefficient
(\MThreeAMake{} pp) exceeds the receiving-payment coefficient (\MThreeAReceive{} pp),
contrary to the prediction that incoming, lender-visible payments matter more. In M3b the
receive interaction is essentially zero (\MThreeBReceive{} pp, $p\approx\MThreeBReceiveP$)
while the make interaction is \MThreeBMake{} pp ($p\approx 6\times10^{-5}$). The
informal-borrowing placebo (M4) also fails: the interaction for borrowing from family or
friends is \textbf{positive} (\InformalBetaThree{} pp, $p\approx\InformalBetaThreeP$;
Benjamini--Hochberg $q\approx\InformalBetaThreeQ$; \InformalN{} adults), where the
substitution reading predicts it to be approximately zero (Table~\ref{tab:secondary}).
Taken together, the secondary results give no support to a credit-information-substitution
interpretation even though some of the individual coefficients are statistically
distinguishable from zero.

% GENERATED by code/study_05_crosscountry/tables/build_tables.py -- DO NOT EDIT BY HAND
\begin{table}[htbp]
\centering
\caption{Mechanism split: receiving versus making digital payments.}
\label{tab:mechanism}

\begin{adjustbox}{max width=\textwidth}
\begin{tabular}{l r c l}
\toprule
 & Coefficient & (CR1 SE) & $p$ \\
\midrule
Receiving a digital payment (M3a) & 4.080*** & (0.578) & $<$0.0001 \\
Making a digital payment (M3a) & 5.172*** & (0.544) & $<$0.0001 \\
\midrule
Receiving $\times$ lower 2019 coverage (M3b) & -0.269 & (0.494) & 0.587 \\
Making $\times$ lower 2019 coverage (M3b) & -2.021*** & (0.480) & $<$0.0001 \\
\midrule
Observations & \multicolumn{3}{c}{90,118} \\
Clusters (economies) & \multicolumn{3}{c}{89} \\
\bottomrule
\end{tabular}
\end{adjustbox}
\smallskip
\begin{minipage}{\linewidth}\footnotesize \emph{Notes.}
Mechanism sub-sample (non-missing receive and make indicators). M3a and M3b
carry the full frozen control set and economy fixed effects. Percentage points;
economy-clustered (CR1) standard errors. The substitution reading predicts the
receiving interaction to be the more negative; it is instead near zero.
\end{minipage}

\end{table}

% GENERATED by code/study_05_crosscountry/tables/build_tables.py -- DO NOT EDIT BY HAND
\begin{table}[htbp]
\centering
\caption{Secondary outcomes and the informal-borrowing placebo.}
\label{tab:secondary}

\begin{adjustbox}{max width=\textwidth}
\begin{tabular}{l r r l r r}
\toprule
Outcome & $\beta_1$ & $\beta_3$ & $p(\beta_3)$ & BH $q(\beta_3)$ & $N$ \\
\midrule
Mobile-loan-app use & 7.375*** & 1.097* & 0.086 & 0.086 & 100,452 \\
Emergency loan from a formal source & 0.718** & -0.883** & 0.022 & 0.029 & 97,540 \\
Informal borrowing (placebo) & 10.662*** & 1.736*** & 0.005 & 0.010 & 100,428 \\
\bottomrule
\end{tabular}
\end{adjustbox}
\smallskip
\begin{minipage}{\linewidth}\footnotesize \emph{Notes.}
Each row is one M2-form regression on the named outcome, full control set and
economy fixed effects, percentage points, economy-clustered (CR1). Benjamini--Hochberg
$q$-values are computed over the frozen nine-hypothesis secondary family. The
substitution reading predicts the informal-borrowing $\beta_3$ near zero; it is
positive and survives the correction.
\end{minipage}

\end{table}


\subsection{Independent reproduction}

The primary and robustness results were produced by one implementation and then reproduced
by a separate cold implementation written from the frozen contracts alone. Of 44 compared
quantities, 41 match to an absolute difference below $5\times10^{-13}$ (largest
$4.8\times10^{-13}$), including M1, the M2 interaction, M2b, the mechanism split, every
secondary outcome, the full robustness grid, and every adversarial check. The three that
differ belong to two non-primary diagnostics --- the M2c joint-test statistic form, and the
logit and probit interaction marginal effects --- and differ in magnitude but not in sign or
conclusion. The reproduction verdict is PASS. The post-registration China/Saudi-Arabia
integrity sensitivity was likewise reproduced by a separate implementation to an absolute
difference of $5\times10^{-14}$.

\section{Vietnam in the cross-country distribution}

This section is descriptive; it makes no claim about a Vietnam-specific effect. Vietnam's
2019 credit-information coverage is \VnmCoverage\% (public registry 59.4\%, private bureau
20.6\%), above the \NEconomies{}-economy median of \SampleMedianCoverage\% --- Vietnam is a
\emph{higher-coverage} economy on this measure, with a maximum-score
depth-of-credit-information index of 8. Evaluated at Vietnam's value of the moderator, the
fitted M2 slope of formal borrowing on digital-payment activity is approximately
$+\VnmSlope$ pp (95\% CI \VnmSlopeCIlo--\VnmSlopeCIhi{} pp; delta method). This is the
cross-country pattern read at Vietnam's 2019 coverage level, not an estimate of what digital
payments do in Vietnam. Both the point estimate and the interval are now persisted
(\texttt{CLAUDE-S5R-M7-VNM-CI-20260910-001}); the point estimate matches the primary run to
twelve decimal places, and the interval is a descriptive diagnostic (M7) outside the frozen
quantity set.

\begin{figure}[htbp]
\centering
\includegraphics[width=0.8\linewidth]{F3_implied_association}
\caption{M2-implied digital-payment/formal-borrowing association across the range of the
2019 coverage moderator (higher $x$ = thinner coverage), with a pointwise 95\% band.
Vietnam's position is marked. The slope of this line is the pre-registered interaction
$\beta_3=\MTwoBetaThree$ pp per SD.}
\label{fig:implied}
\end{figure}

\section{Discussion}

\paragraph{Why might the association be weaker, not stronger, where coverage is thin?} Three
explanations are consistent with the evidence and with the design's limits.

\emph{A coverage floor.} In the \ThinCoverageEconomies{} economies with 2019 coverage below
10\%, credit-information systems are close to absent for everyone; there is little
cross-sectional room for a digital-payment record to ``substitute'' at the margin, and
formal borrowing in these economies is low and concentrated among a small, selected group.
Dropping the zero-coverage economies and the lowest coverage quartile does not remove the
negative interaction, which argues against this being the whole story, but a non-linear
coverage relationship remains plausible.

\emph{The account-access channel.} Where formal financial infrastructure is thin, holding an
account and using digital payments may principally mark inclusion in the formal system at
all, rather than an incremental, lender-legible signal on top of an existing credit file.
The account-holders-only restriction shrinks the interaction toward zero and renders it
imprecise, which is consistent with composition --- including non-account holders ---
accounting for part of the gradient. The fully interacted-control check also shows partial
attenuation. These clues make the account-access channel a plausible interpretation, but
neither identifies it: the restriction is imprecise and conditions on a possible collider,
while coefficient attenuation can have several causes. We therefore treat this as a
hypothesis for future data, not the paper's mechanism finding.

\emph{Supply-side heterogeneity.} Lenders in thin-coverage economies may not have the
systems to act on a digital-payment history even when it exists; the firm-level result may
reflect relationship-lending technologies that individual retail lending in these settings
lacks. We cannot test this with Findex.

\paragraph{Reconciling with \citet{GalileaFaraziMare2026}.} The unit differs (firms versus
individuals), the outcome differs (self-reported credit constraints versus realised formal
borrowing), the moderator differs (depth index versus coverage), and the lending
technologies differ. A contrary individual-level cross-section does not contradict their
firm-level finding; it bounds its external validity.

\paragraph{Policy boundary.} The rejected moderation prediction does not imply that digital
payments fail to improve access, that credit reporting is unimportant, or that payment data
cannot aid underwriting. It says only that this 2024 individual cross-section does not show
a larger conditional payment--borrowing association in economies with thinner archived 2019
coverage.

\paragraph{Residual limitations that must survive typesetting.} The following limitations
are carried verbatim from the independent adversarial review and are not weakened anywhere
in this paper:
\begin{itemize}
\item The design is a single cross-section; reverse causality and selection remain.
\item The cross-level interaction is identified by \NEconomies{} economy values of an
archived 2019 moderator, not by \NAdults{} independent country observations.
\item The Doing Business checks establish agreement with archived profiles/live World Bank
data, not validation against a fully independent measurement system. General non-classical
measurement error remains possible.
\item The depth and legal-rights proxies do not significantly reproduce the coverage result.
\item The C5 fully interacted model is post-registration. Its coefficient is
dual-implementation; its within-$R^2$ diagnostic is persisted but single-implementation.
\item The MDE is a retrospective precision summary using the observed clustered SE, not an
ex ante design guarantee or proof of a negative structural effect.
\item The account-access channel is a plausible hypothesis only; it is not identified.
\item \texttt{search\_adequacy} remains \texttt{PARTIAL}, so the paper cannot claim to be
the first study.
\end{itemize}
Additional design-level caveats: \texttt{account\_fin} is conditioned on but is a possible
common consequence, so the account-holders-only check is a sensitivity analysis, not a fix;
the private-credit robustness control has heterogeneous vintages; and the Global Findex
provenance is verified-reconstructed, not original-download verified. None of these is
resolved by the design, and all are why the language stays associational.

\section{Conclusion}

In a pre-registered individual-level test across \NEconomies{} economies, digital-payment
activity is positively associated with formal borrowing, but the
credit-information-substitution prediction --- that the association is larger where
credit-information coverage is thinner --- is not supported. For the registered coverage
moderator, the interaction has the opposite sign across weighting, sample and control-set
variants and all leave-one-economy-out fits; the two alternative institutional proxies are
imprecise. The registered mechanism and placebo tests give no support to the substitution
reading. We read this as an external-population check that does not reproduce the firm-level
pattern at the individual level in this cross-section. Basic formal-system access is one
plausible explanation, but the present design cannot identify it. The value of the exercise
is in its transparency: a frozen plan, an independently reproduced result, and a contrary
finding reported as such.

\section*{Reproducibility statement}

Raw inputs, code and results are under \texttt{projects/study\_05\_findex\_crosscountry/}.
The frozen pre-registration and hypothesis file were hash-committed before estimation.
Primary estimation and independent reproduction each have a fixed seed and a recorded
input/output SHA-256; every table and in-text statistic in this manuscript is generated by
\url{code/study_05_crosscountry/tables/build_tables.py} from those immutable result
files. Input SHA-256: Findex microdata \texttt{ca307a0c\ldots f90fce}; 2019
credit-information snapshot \texttt{ef0c67f9\ldots bd06c}.
